\chapter{L'évasion}
\label{chap:evasion}

Khaer Ghal attendit la nuit.

Depuis que son père lui avait annoncé que la famille retournerait dans l'arène
le lendemain, il avait passé des heures à chercher une autre solution. Il
avait déjà essayé de discuter avec lui, puis de s'interposer lorsque les
créatures avaient été lâchées. La première tentative n'avait rien changé ; la
seconde n'avait offert à Kamatsu et à ses parents qu'un jour supplémentaire.

Cette fois, il ne pouvait plus se contenter de retarder ce que son père avait
décidé.

La punition reçue sur la place du village rendait pourtant chacun de ses
mouvements douloureux. Les coups infligés par son père s'ajoutaient aux
blessures de l'arène et, lorsqu'il se releva de sa couche, son corps lui
rappela immédiatement qu'il aurait probablement dû rester allongé.

Il se leva malgré tout.

Le village était calme. Quelques foyers achevaient de se consumer devant les
habitations, mais la plupart des Orcs dormaient depuis longtemps. Seuls les
gardes poursuivaient leurs rondes le long de la palissade, près de l'entrée et
autour de l'enclos des prisonniers.

Khaer Ghal connaissait leurs habitudes. Il avait grandi parmi eux, joué entre
ces maisons et observé pendant des années les guerriers surveiller le village.
Il savait à quel moment certaines parties de la palissade disparaissaient de
leur champ de vision, quel guerrier avait l'habitude de s'attarder quelques
minutes lorsqu'il relevait son compagnon et quels passages permettaient de
traverser les habitations sans croiser les zones encore éclairées.

Tout ce qu'il avait appris pour connaître et protéger son village allait
maintenant lui servir à agir contre lui.

Il commença par préparer ce dont les fugitifs auraient besoin. Quelques
morceaux de viande séchée, du pain, une outre d'eau et une couverture
disparurent dans un sac. Il hésita devant les réserves, calcula grossièrement
ce que trois personnes pourraient consommer pendant plusieurs jours, puis
ajouta encore de la nourriture.

Il leur fallait aussi de quoi se défendre. Parmi les armes prises aux voyageurs
lors de l'attaque de la caravane, Khaer Ghal récupéra deux lames et une petite
hache. Il reconnut certaines d'entre elles : quelques jours auparavant, il les
avait vues ramassées sur la route par les guerriers de son village.

Cette nuit, il les rendrait à leurs propriétaires.

Restait la palissade.

Khaer Ghal avait choisi son passage avant même la tombée de la nuit. Sur le
côté du village opposé à l'entrée principale, une portion de l'enceinte avait
été endommagée quelques mois auparavant par la chute d'un arbre. Les grands
pieux verticaux avaient été remis en place, mais la réparation n'était pas
aussi solide que le reste de la fortification : plusieurs rondins plus courts
et des pièces de renfort avaient été ajoutés entre les troncs d'origine, puis
maintenus par de lourdes attaches de corde et des chevilles de bois.

Il avait observé les adultes effectuer cette réparation. Il savait comment
elle tenait et, surtout, comment la défaire.

Il ne restait plus qu'à attendre le bon moment.

\scenebreak

Lorsque le garde passa devant l'enclos et disparut derrière une habitation,
Khaer Ghal traversa rapidement l'espace découvert. Kamatsu dormait contre sa
mère ; son père reposait non loin d'eux, suffisamment près pour que le moindre
mouvement de son fils le réveille probablement.

Khaer Ghal s'accroupit près de la clôture.

— Kamatsu.

Le garçon ne réagit pas.

— Kamatsu.

Cette fois, il remua. Ses yeux s'ouvrirent et il lui fallut quelques secondes
pour reconnaître Khaer Ghal dans l'obscurité avant de se redresser
brusquement.

— Qu'est-ce que tu fais ?

— Moins fort.

Les parents de Kamatsu s'étaient réveillés à leur tour. Son père se leva
immédiatement et s'approcha de la clôture avec méfiance tandis que Khaer Ghal
posait le sac au sol.

— Vous devez partir.

L'homme le fixa.

— Quoi ?

— Maintenant.

Khaer Ghal sortit la clé qu'il avait réussi à récupérer et ouvrit l'enclos.
Personne ne bougea. Il poussa légèrement la porte.

— Demain, vous retournez dans l'arène.

La mère de Kamatsu échangea un regard avec son mari.

— On le sait.

— Alors partez.

Khaer Ghal leur tendit le sac avant de poser les armes devant eux. Le père de
Kamatsu regarda successivement les lames, la porte ouverte et le jeune Orc.

— Pourquoi ?

Khaer Ghal fronça les sourcils. Compte tenu du temps dont ils disposaient, la
question lui semblait particulièrement mal choisie.

— Parce que si vous restez, vous allez mourir.

L'homme continua de le regarder.

— Et toi ?

Khaer Ghal hésita. Il avait pensé à l'enclos, aux gardes, à la nourriture, aux
armes, à la palissade et au chemin qu'ils devraient emprunter dans la forêt.
Il avait beaucoup moins réfléchi à ce qui lui arriverait après.

— Moi, j'habite ici.

Kamatsu le regarda avec incrédulité.

— Ils vont savoir que c'est toi.

— Je sais.

— Et ton père ?

Khaer Ghal détourna brièvement les yeux.

— Dépêchez-vous.

Le père de Kamatsu comprit qu'il n'obtiendrait rien de plus. Il ramassa les
armes, en tendit une à sa femme et passa le sac sur son épaule.

Cette fois, personne ne discuta.

\scenebreak

Khaer Ghal les guida à travers le village en empruntant les passages qu'il
avait repérés. Il avançait de quelques mètres, s'arrêtait pour écouter puis
leur faisait signe de le rejoindre. Derrière lui, les trois humains
progressaient aussi silencieusement qu'ils le pouvaient, mais leur manque
d'habitude lui paraissait assourdissant. Une botte racla la terre, un vêtement
accrocha un morceau de bois, une pierre roula sous un pied ; à chaque bruit,
Khaer Ghal s'immobilisait en retenant son souffle.

À deux reprises, ils durent se cacher. Une première patrouille passa
suffisamment près pour qu'il entende la respiration de l'un des guerriers. Plus
loin, une porte s'ouvrit brusquement et les quatre fugitifs se plaquèrent
derrière une habitation tandis qu'un Orc traversait le chemin avant de
disparaître entre les bâtiments.

Personne ne donna l'alerte.

Ils atteignirent enfin la palissade.

La portion réparée était exactement telle que Khaer Ghal l'avait laissée dans
son souvenir. Entre les énormes troncs formant l'enceinte, plusieurs pièces de
bois plus courtes renforçaient la partie reconstruite. Il s'agenouilla devant
les premières attaches, sortit son couteau et entreprit de couper la corde.

Elle résista.

— Aide-moi.

Le père de Kamatsu s'accroupit auprès de lui. Ensemble, ils défirent une
première fixation, puis une seconde. Une traverse fut retirée et déposée
lentement au sol pour éviter qu'elle ne tombe avec fracas. Une autre suivit.

Le travail prit beaucoup plus de temps que Khaer Ghal l'avait espéré. Peu à
peu pourtant, la réparation céda. Ils dégagèrent plusieurs rondins et firent
basculer l'un des pieux les plus minces vers l'intérieur. Une ouverture apparut
dans la fortification, d'abord étroite, puis suffisamment large pour laisser
voir les arbres au-delà.

Le père de Kamatsu tira encore sur une pièce de bois.

Elle céda.

Une véritable brèche s'ouvrait désormais devant eux. Elle n'occupait qu'une
portion limitée de la palissade, mais plusieurs éléments avaient disparu ou
basculé, créant un passage irrégulier assez large pour qu'un adulte puisse le
franchir rapidement. De chaque côté, les lourds troncs de l'enceinte restaient
debout et donnaient à l'ouverture l'apparence d'une blessure grossièrement
arrachée dans la réparation plutôt que celle d'une fortification effondrée.

L'air froid de la forêt passa à travers.

Pendant quelques secondes, personne ne bougea. De l'autre côté, il n'y avait
plus de village, seulement les arbres et la nuit. Khaer Ghal regarda
l'ouverture, puis Kamatsu.

— Continuez vers l'est jusqu'au ruisseau, murmura-t-il. Ensuite, suivez l'eau
vers le sud. Ne prenez pas les chemins. Ils vous chercheront d'abord dessus.

Le père de Kamatsu hocha la tête. Sa femme franchit la brèche la première et
Kamatsu s'engagea derrière elle, avant de se retourner.

— Viens avec nous.

Khaer Ghal le regarda comme s'il venait de proposer quelque chose d'absurde.

— C'est mon village.

— Après ce que tu viens de faire ?

Khaer Ghal ouvrit la bouche, mais aucun mot ne vint.

Un bruit derrière lui mit fin à la conversation.

Des pas.

Beaucoup trop de pas.

Khaer Ghal se retourna. Plusieurs silhouettes venaient d'apparaître entre les
habitations. Des guerriers orcs avançaient dans leur direction, suffisamment
nombreux et suffisamment espacés pour leur interdire tout retour vers le cœur
du village.

Au milieu d'eux marchait son père.

Khaer Ghal s'immobilisa.

%\illustration
%  {0600.png}
%  {La fuite interrompue}
%  {evasion}

Pendant un instant, personne ne parla. Le regard du chef passa de son fils aux
armes rendues aux prisonniers, puis à la large ouverture pratiquée dans la
palissade. Il n'avait besoin d'aucune explication.

Khaer Ghal non plus.

Ils avaient été découverts.

Derrière lui, pourtant, Kamatsu et sa mère se trouvaient déjà dans la forêt.
Le père de Kamatsu était encore à ses côtés, juste devant la brèche. Il restait
une possibilité.

Khaer Ghal se retourna vers eux.

— Courez.

Le père de Kamatsu hésita.

— Courez !

Sa femme saisit Kamatsu par le bras et l'entraîna entre les arbres. Son père
fit un pas vers l'ouverture, puis regarda les guerriers qui approchaient.

Il ne les suivit pas.

Il resserra les doigts autour de son arme et se plaça devant la brèche. Khaer
Ghal comprit immédiatement.

— Va avec eux.

L'homme secoua la tête.

Il occupait presque toute la largeur du passage. La brèche qui leur avait
permis de sortir rapidement devenait maintenant un étranglement naturel : pour
poursuivre sa femme et son fils, les Orcs devraient d'abord passer près de lui.

Les premiers guerriers arrivèrent et Khaer Ghal tira sa dague.

Il savait ce qu'on attendait de lui. Le fils du chef qu'on avait formé pendant
toutes ces années aurait dû obéir, laisser les guerriers reprendre les
prisonniers et accepter la punition qui lui serait destinée.

Mais il avait déjà fait son choix.

Il avait ouvert l'enclos parce qu'il l'avait décidé. Il avait rendu leurs armes
aux prisonniers parce qu'il l'avait décidé. Il avait ouvert cette palissade
parce qu'il refusait de laisser son père choisir qui devait mourir le
lendemain.

Pour la première fois, Khaer Ghal ne cherchait plus à obtenir la permission de
faire ce qu'il estimait juste.

— Partez, répéta-t-il.

Puis les guerriers furent sur eux.

\scenebreak

Khaer Ghal para le premier coup avec sa dague, mais la force de l'impact lui
arracha presque l'arme des doigts. Il tenta de riposter et reçut aussitôt un
violent revers qui l'envoya rouler dans la poussière. Son corps, déjà meurtri
par l'arène et la punition de la veille, n'avait plus grand-chose à offrir.

Il voulut se relever, mais une silhouette s'arrêta devant lui.

Son père.

Khaer Ghal resta à genoux, la dague toujours serrée dans sa main. Il avait été
capable de désobéir à cet homme, de libérer ses prisonniers et même de lever
une arme contre les guerriers venus les reprendre.

Mais son père était autre chose.

C'était l'homme qui l'avait élevé, celui dont il avait recherché
l'approbation depuis qu'il était suffisamment âgé pour comprendre ce qu'elle
signifiait. Malgré tout ce qui venait de se passer, malgré les coups de la
veille et la colère qui brûlait encore en lui, Khaer Ghal ne parvint pas à
lever sa lame contre lui.

Sa main descendit lentement et la dague tomba dans la poussière.

Deux guerriers le saisirent aussitôt.

Près de la palissade, le père de Kamatsu défendait toujours la brèche. Sa
position empêchait plusieurs Orcs de l'attaquer ensemble et il parvint à
repousser un premier adversaire avant de faire reculer le suivant. Chaque
seconde gagnée permettait à sa femme et à son fils de s'enfoncer davantage
dans la forêt.

Mais il n'était pas un guerrier.

Un premier coup l'atteignit et le fit chanceler. Il retrouva son équilibre,
leva de nouveau son arme et tenta de tenir encore quelques instants.

Le coup suivant passa sous sa garde.

Il tomba.

— Non !

La voix de Kamatsu résonna au-delà de la palissade.

Son père ne se releva pas.

Khaer Ghal se débattit entre les deux Orcs qui le retenaient.

— Laissez-les partir !

Personne ne l'écouta. Deux guerriers franchirent la brèche et disparurent entre
les arbres. Un troisième s'arrêta près de l'ouverture et porta la main à son
arc.

Khaer Ghal comprit immédiatement.

— Non !

L'Orc saisit une flèche et arma. À travers la brèche, entre les troncs, Khaer
Ghal distinguait encore la mère de Kamatsu qui entraînait son fils avec elle.

La corde claqua.

La flèche disparut dans l'obscurité.

La femme s'effondra.

— Maman !

Kamatsu s'arrêta. Il aurait encore pu courir : la forêt s'étendait devant lui
et les deux guerriers lancés à leur poursuite n'étaient pas encore assez
proches pour le rattraper.

Mais sa mère était au sol.

Kamatsu revint vers elle et tomba à genoux à ses côtés.

— Continue ! hurla Khaer Ghal.

Son ami ne l'entendit pas, ou ne voulut pas l'entendre. Il resta auprès de sa
mère jusqu'à ce que les deux Orcs le rejoignent.

Lorsqu'ils le saisirent, il ne tenta même pas de fuir.

\scenebreak

Les deux garçons furent ramenés au village. Cette fois, personne ne parla
pendant le trajet. Kamatsu avançait entre deux guerriers, le regard fixé devant
lui sans paraître réellement voir ce qui l'entourait. Khaer Ghal marchait
quelques pas derrière, les poignets maintenus dans son dos.

Il aurait voulu trouver quelque chose à dire.

Aucun mot ne pouvait rendre à Kamatsu ses parents.

Lorsqu'ils atteignirent la place centrale, son père les attendait. Le chef
regarda d'abord Kamatsu, puis son fils.

— Je t'avais ordonné de rester à ta place.

Khaer Ghal soutint son regard.

— Ils allaient mourir.

Son père observa un instant le garçon que les guerriers maintenaient à quelques
pas de là.

— Et maintenant, ses parents sont morts.

Khaer Ghal serra les mâchoires.

— Tu as désobéi à ton chef. Tu as libéré des prisonniers. Tu leur as rendu des
armes prises par le village. Tu as détruit une partie de notre palissade et tu
as combattu ceux qui tentaient de les reprendre.

Chaque accusation tombait avec le calme d'un jugement déjà rendu.

Son père désigna Kamatsu.

— Tout ça pour lui.

Khaer Ghal regarda Kamatsu.

— Oui.

Le mot sortit sans hésitation.

Son père resta silencieux quelques secondes, comme s'il attendait encore autre
chose : une excuse, un regret, peut-être simplement que son fils baisse les
yeux.

Khaer Ghal n'en fit rien.

Le chef s'approcha.

— Puisque tu tiens tant à lui, tu vas finir comme lui.

Il fit signe aux guerriers.

— Enfermez-les ensemble.

Khaer Ghal eut un mouvement de surprise. Cette fois, les mains qui le
retenaient ne le conduisirent ni vers sa maison ni vers la place où il avait
été puni la veille. Elles le poussèrent vers l'enclos des prisonniers.

La porte s'ouvrit. Kamatsu fut conduit à l'intérieur le premier, puis Khaer
Ghal le suivit avant que le verrou ne retombe derrière eux.

Pendant un long moment, aucun des deux garçons ne parla. Kamatsu finit par
s'asseoir contre la clôture, les genoux ramenés contre lui, le regard fixé au
sol. Khaer Ghal resta debout quelques instants avant de venir s'installer à
quelques pas de lui.

Il aurait voulu s'excuser. Il avait déjà essayé devant l'arène, lorsqu'il
n'avait pas réussi à empêcher le combat. Cette fois, le mot lui parut
dérisoire.

Alors il resta simplement là.

Pour la première fois de sa vie, Khaer Ghal regardait son village depuis
l'intérieur de l'enclos des prisonniers.